Tras posicionar una herramienta en una bancada, su pequeño desplazamiento libre se modela por
\[
x''+4x=0,
\]
con \(t\) en segundos y \(x\) en milímetros; el coeficiente cuatro tiene unidades \(\mathrm{s}^{-2}\). Los errores de arranque sólo garantizan
\[
-1\le x(0)\le1\ \mathrm{mm},\qquad
-2\le x'(0)\le2\ \mathrm{mm/s}.
\]
Cualquier pareja en ese rectángulo es admisible; no se asignan probabilidades.

Determine cotas superior e inferior ajustadas para \(x(t)\) en cada instante y dibújelas a mano en \(0\le t\le\pi/2\) s. Explique oralmente si una sola trayectoria alcanza toda la envolvente superior. Decida si esos datos garantizan la holgura \(|x(t)|\le1.30\,\mathrm{mm}\) durante ese intervalo, distinguiendo una garantía para todos los arranques de lo que puede ocurrir con uno particular.
